% ============================================================
% Section 2: Background
% ============================================================
% Target length: ~0.75 column (400–500 words)
% Write: Month 4, build incrementally
% Structure: capability systems → agentic browsers →
%            audit logs → why combining them is novel
% ============================================================

\section{Background}
\label{sec:background}

\subsection{Capability-Based Security}

Capability systems restrict access to resources by requiring
unforgeable tokens (capabilities) that confer specific rights.
Capsicum~\cite{watson2010capsicum} extended POSIX with
capability-mode processes and capability file descriptors,
preventing applications from accessing resources beyond
explicitly granted capabilities.
COWL~\cite{stefan2014cowl} brought label-based information
flow control to web browsers, restricting cross-origin
communication for browser workers.
RLBox~\cite{narayan2020rlbox} sandboxed third-party libraries
in Firefox using WebAssembly, demonstrating capability isolation
within a production browser context.
None of these systems addressed the agentic context where
an AI principal acts with user-level privileges.

\subsection{Agentic Browsers}

\todo{Write when product landscape is fully confirmed.
      Cover: Perplexity Comet, ChatGPT Atlas, Dia (Atlassian),
      Brave Leo, Chrome/Edge Copilot integrations.
      Key point: all grant agents ambient authority with no
      formal capability model.}

\subsection{Prompt Injection}

Prompt injection attacks embed malicious instructions in content
that AI agents process, causing them to perform unauthorized
actions~\cite{greshake2023not}.
Indirect injection---where instructions are hidden in web content
rather than the user's direct input---is particularly dangerous
in browser contexts~\cite{nasr2025}.
\todo{Add SecureWebArena, WASP, BrowseSafe citations when confirmed.}

\subsection{Cryptographic Audit Logs}

Certificate Transparency (CT)~\cite{laurie2013rfc6962} introduced
append-only Merkle logs for TLS certificate issuance, enabling
public auditability of certificate authority behavior.
Sigstore/Rekor~\cite{sigstore_rekor} extended this model to
software supply chain verification.
\ferrite{} applies the same principle to browser agent actions:
every capability grant, denial, and exercise is recorded in a
hash-chained log (upgrading to a Merkle tree), enabling
post-hoc verification that every agent action was authorized.

\subsection{The Gap}

\todo{One paragraph synthesising why no existing work combines
      these three areas. This is the core research gap claim.
      Reference the table from PROJECT\_REFERENCE.md.}
